\section{Binomial phase and the optimal number of zero coordinates}
\label{sec:iid-phase}

This section concerns the binomial optima $B_m$ and their padded envelope
$P_N:=\max_{1\le m\le N}B_m$. It determines the eventual record pattern
within this family. Theorem~\ref{thm:common} transfers these conclusions
to common-diameter maximizers; no reduction from arbitrary independent
laws is assumed here.

Put $s=\sqrt{10}$ and write
\begin{equation}\label{eq:phase-constants}
\begin{gathered}
 p=(4-s)/2,\quad q=1-p,\quad t=s-3,\quad u=p-t,
 \quad c_*=(20+s)/72,\\
 Q_*=\frac{\phz(867+287s)}{25920},\qquad
 \alpha=\frac{\phz(60+19s)}{30},\qquad
 b=\frac{1010-397s}{12960},\\
 \delta=c_*-b=\frac{2590+577s}{12960},\qquad
 \eta=Q_*/\alpha=\frac{2510-747s}{8640},\qquad
 \omega=\frac{7773-2447s}{8640}.
\end{gathered}
\end{equation}
The letters $p,q,t$ in this section without an argument denote these
fixed limiting values. For a variable Bernoulli parameter $v$, use
$a(v)=v(1-v)$, $r(v)=v^2+(1-v)^2$ and $\tau(v)=1-2v$.
For a real number $x$, write $\langle x\rangle=x-\operatorname{round}(x)$,
choosing either convention at half-integers.

\begin{theorem}[Eventual binomial padding]\label{thm:iid-padding}
At every sufficiently large horizon $N$, $P_N$ has a unique maximizing
core size $m_N$, and $j_N=N-m_N$ belongs to $\{0,1,2,3,4\}$.
The limiting frequencies of these five zero counts are respectively
\[
 p-\omega,\quad p-\omega,\quad t,\quad\omega,\quad\omega.
\]
The gaps between successive record core sizes eventually belong exactly
to $\{2,3,5\}$: no other gap occurs sufficiently late, and each of
these three occurs infinitely often. No effective first horizon is asserted.
\end{theorem}

\subsection{A uniform expansion at a fixed integer threshold}

Let $F_{m,v}$ be the distribution function of $\Bin(m,v)$. Fix an integer
$k$ and set $c=mv-k$, $d=(ma(v))^{-1/2}$. The positive closed-threshold
discrepancy is $D_{m,k}(v)=F_{m,v}(k)-\Phi(-cd)$.

\begin{lemma}\label{lem:phase-uniform}
On every compact $J\subset(0,1)$ and bounded range $|c|\le M$, uniformly
for $v=(k+c)/m\in J$,
\begin{equation}\label{eq:phase-cdf}
 D_{m,k}(v)=\phz\bigl[d(1/2+\tau/6)+d^3C(\tau,-c)
                         +d^5T_5(\tau,-c)\bigr]+O(d^7),
\end{equation}
with the same order for the derivative of the remainder in $c$, fixing
$m,k$. The expansion extends to any fixed order and any fixed finite
number of these branch derivatives. Here
\begin{align*}
C(\tau,\beta)={}&\frac{-1080(1+\tau)\beta^2
 -(90\tau^2+1080\tau+630)\beta
 +\tau^3-45\tau^2-189\tau-135}{4320},\\
T_5(\tau,\beta)={}&\frac{3+5\tau}{48}\beta^4
 +\frac{27+60\tau+25\tau^2}{288}\beta^3
 +\frac{15+31\tau+25\tau^2+\tau^3}{192}\beta^2\\
&+\frac{549+1320\tau+1030\tau^2+120\tau^3+5\tau^4}{23040}\beta\\
 &+\frac{189+927\tau+126\tau^2+234\tau^3+21\tau^4-25\tau^5}{193536}.
\end{align*}
\end{lemma}

\begin{proof}
The Fourier integral for a binomial mass is on $[-\pi,\pi]$. Its modulus
satisfies
\[
 |1-v+ve^{i\xi}|^m
 =(1-4a(v)\sin^2(\xi/2))^{m/2}.
\]
After $z=\xi\sqrt{ma(v)}$, the inequality
$\sin(|\xi|/2)\ge|\xi|/\pi$ bounds this by $e^{-2z^2/\pi^2}$
throughout the integration interval. On $J$, the centered logarithm has
a Taylor expansion of any prescribed order with uniformly bounded
coefficients and parameter derivatives. Split the integral at
$|z|=m^\varepsilon$, with $\varepsilon$ sufficiently small for that
order. Polynomial multiples of the Gaussian bound control the inner
Taylor remainder; the outer integral is smaller than any prescribed
power of $m^{-1}$. Taking sufficiently many terms gives a Hermite
mass expansion with uniform absolute error $O(d^{12})$, or any higher
fixed power if needed, at every integer mass index.

This statement also holds for branch derivatives. Before centering,
\[
 \frac{\partial}{\partial c}(1-v+ve^{i\xi})^m
 =(e^{i\xi}-1)(1-v+ve^{i\xi})^{m-1},
 \qquad v=(k+c)/m.
\]
Further fixed derivatives have the same Gaussian control after increasing
the retained Taylor order. The standardized mass index and spacing have
bounded branch derivatives of the needed orders on the stated ranges.

Summing at most $m+1=O(d^{-2})$ masses turns the displayed mass error
into $O(d^{10})$, including its branch derivatives. The tails beyond
the true binomial support are exponentially small. Euler--Maclaurin
summation may therefore be applied to the infinite lower lattice of the
Gaussian-polynomial terms. Its order-eight remainder is bounded by a
constant times $d^8\|f^{(8)}\|_1$ for each such term. In the integral
periodic-Bernoulli remainder, differentiating the phase gives bounded
factors of the form $1+d|z|$; Gaussian derivative moments remain
integrable. Increasing the summation order gives the corresponding
assertion for further branch derivatives.

The coefficient of $d^j$ in the density has parity $j$. Integration,
the half-cell term, and the Euler--Maclaurin derivative terms then leave
only odd powers $d,d^3,d^5,\ldots$ when evaluated at $-cd$.
The terms through $d^5$ give \eqref{eq:phase-cdf}. Their coefficients
follow by using the Bernoulli cumulant ratios
\[
 \tau,\quad1-6a,\quad\tau(1-12a),\quad
 1-30a+120a^2,\quad\tau(1-60a+360a^2)
\]
and the lattice corrections $df/2+d^2f'/12-d^4f'''/720$.
The cubic phase coefficient is also the polynomial in Schulz's
Remark~4.3, derived on p.~25 \cite{Schulz2016}.
\end{proof}

Dividing \eqref{eq:phase-cdf} by $r(v)d$ gives the threshold efficiency
\begin{equation}\label{eq:phase-efficiency}
 R_{m,+,k}(v)=E(v)-\frac{Q(v,c)}m-\frac{H(v,c)}{m^2}+O(m^{-3}),
\end{equation}
where
\[
 E(v)=\phz\frac{3+\tau(v)}{3(1+\tau(v)^2)},\quad
 Q(v,c)=-\frac{\phz C(\tau(v),-c)}{r(v)a(v)},\quad
 H(v,c)=-\frac{\phz T_5(\tau(v),-c)}{r(v)a(v)^2}.
\]
The phase parabola $Q(v,c)$ has vertex
\[
 c_0(v)=\frac{\tau(v)^2+12\tau(v)+7}{24(1+\tau(v))}
\]
and minimum
\begin{equation}\label{eq:phase-qmin}
 Q_{\min}(v)=\frac{\phz\{345+72\tau-498\tau^2-8\tau^3-23\tau^4\}}
 {4320(1+\tau)(1-\tau^4)}.
\end{equation}

\subsection{Global binomial localization and the padding envelope}

Schulz's Theorem~1 gives $B_m<\ce$ for every finite $m$.
His Proposition~6.1 bounds the full Kolmogorov distance on
$v\in[1/6,5/6]$ by the four absolute closed/left central-threshold errors
at $\lfloor mv\rfloor,\lceil mv\rceil$ and
\[
 \gamma\frac{\phz(3+|1-2v|)}{6\sqrt{mv(1-v)}},
 \qquad \gamma=\frac{2}{\sqrt{2\pi}\ce}-1=12t-1<1.
\]
His estimate on the two complementary parameter ranges is
$R_m(v)\le .3031(1+.646/\beta(v))<.4045$,
where $\beta(v)=r(v)/\sqrt{a(v)}$
\cite[pp.~29--31]{Schulz2016}.
Each $B_m$ is attained: the ratio is the maximum of finitely many
continuous threshold expressions and tends to zero as $v\to0,1$.

By reflection restrict $v\le1/2$. Central expansions and the source
bound show uniform convergence of the optimized threshold value to $E(v)$
on the retained compact parameter interval. Its unique maximum is
$\ce$ at $v=p$. Thus every maximizing parameter tends to $p$.
Near $p$, the extra $\gamma$ term and left-limit thresholds are uniformly
smaller than the positive closed central threshold: the latter has leading
coefficient $\phz(1/2+\tau/6)/r$, while the left-limit coefficient is
$\phz(1/2-\tau/6)/r$. Consequently the full winner is eventually
a positive closed floor/ceiling threshold, with $|c|\le1$.

Taking $v=(k+c_*)/m$ with $k$ nearest $mp-c_*$ gives
$\ce-Q_*/m+O(m^{-2})$. Since
$E(v)=\ce-\alpha(v-p)^2+O((v-p)^3)$, comparison with
\eqref{eq:phase-efficiency} and $Q\ge Q_{\min}$ yields
\[
 (v-p)^2\le C|v-p|/m+C/m^2,
 \qquad v-p=O(m^{-1}).
\]
The positive quadratic phase penalty then gives
$c-c_0(v)=O(m^{-1/2})$. Since $0<c_*<1$, the branch is eventually
interior. Differentiating \eqref{eq:phase-efficiency} along
$v=(k+c)/m$ and using $E'(v)=O(m^{-1})$ gives $Q_c=O(m^{-1})$,
so $c-c_0(v)=O(m^{-1})$ and
$v=(k+c_*)/m+O(m^{-2})$.

Set $e=k+c_*-mp$. Taylor expansion and completion of the square give
\begin{equation}\label{eq:phase-bm}
 B_m=\ce-\frac{Q_*}m-\frac{H_*+
 \alpha\operatorname{dist}(mp-\delta,\Z)^2}{m^2}+O(m^{-3}).
\end{equation}
Indeed the second-order branch correction is
$\alpha e^2+Q_p e+H_0$, where all derivatives are evaluated at
$(p,c_*)$, $H_0=H(p,c_*)$, $b=-Q_p/(2\alpha)$ and
$H_*=H_0-\alpha b^2$. Direct substitution in
\eqref{eq:phase-qmin} gives \eqref{eq:phase-constants}.
The formula remains uniform at the nearest-integer crossing, by taking
the better of the two branch values; the nearest branch itself need
only be selected outside an $O(m^{-1})$ half-integer neighborhood.

Choose the most recent maximizing core $m_N$ in a tie. Strict finite
$B_m<\ce$ and $B_m\to\ce$ imply $m_N\to\infty$.
The leading correction in \eqref{eq:phase-bm}, compared with $B_N$,
then gives $m_N/N\to1$, and its $O(m^{-2})$ bounds give
$j_N=N-m_N=O(1)$. With
$\theta_N=\operatorname{round}(Np-\delta)-(Np-\delta)$, for bounded $j$,
\[
 N^2(B_{N-j}-B_N)=\alpha\{\theta_N^2-
 \operatorname{dist}(\theta_N+jp,\Z)^2-\eta j\}+O(N^{-1}).
\]
The limiting policy minimizes
$f_j(\theta)=\eta j+\operatorname{dist}(\theta+jp,\Z)^2$.
Define
\begin{equation}\label{eq:phase-boundaries}
\begin{aligned}
 T_{13}&=-(p+u)/2+\eta/t,& T_{30}&=-u/2-3\eta/(2u),\\
 T_{02}&=t/2+\eta/t,&T_{24}&=3t/2+\eta/t,\\
 T_{41}&=(q+2t)/2-3\eta/(2u).
\end{aligned}
\end{equation}
In increasing order on $[-1/2,1/2]$ the winning zero counts are
\begin{center}
\begin{tabular}{cc}
\toprule
Phase interval&Zero count\\
\midrule
$[-1/2,T_{13}]$&1\\
$[T_{13},T_{30}]$&3\\
$[T_{30},T_{02}]$&0\\
$[T_{02},T_{24}]$&2\\
$[T_{24},T_{41}]$&4\\
$[T_{41},1/2]$&1\\
\bottomrule
\end{tabular}
\end{center}
The relevant parabola centers are $-p,-u,0,t,2t,q$, with biases
$\eta,3\eta,0,2\eta,4\eta,\eta$. Subtracting neighboring
quadratics gives exactly \eqref{eq:phase-boundaries}. The remaining
copies are above this envelope: their differences are affine, so it
suffices to check the displayed interval endpoints. The successive
active widths are $\omega,p-\omega,t,\omega$, with the two outer
pieces totaling $p-\omega$. Positivity follows, in particular, from
$7773^2-10(2447)^2=541439>0$.

For completeness all $j\ge5$ have a uniform positive price gap. Using
only $j=0,1,2$, the relevant centers are $-p,0,t,q$ and the intersections
are $-p/2-\eta/(2p)$, $t/2+\eta/t$ and
$(q+t)/2-\eta/(2p)$. The largest potentially troublesome intersection
price is at most $2\eta+p^2/4<5\eta$, with gap greater than
$59/10000$, using
\[
 4/25<t<1/6,\qquad5/12<p<21/50,\qquad1/60<\eta<1/55.
\]
The other intersection prices are below $1/16<5\eta$;
the endpoint price $\eta+t^2/4$ is also below $5\eta$.
But $f_j\ge5\eta$ for $j\ge5$. The bounded-$j$ comparison
therefore gives $j_N\le4$ eventually.

Away from $O(N^{-1})$ neighborhoods of the five crossings, the actual
winner has the displayed code. Since $p$ is irrational, its rotation
is equidistributed. Those shrinking neighborhoods have density zero:
their upper density is bounded by the length of any fixed larger
neighborhood. The interval widths prove the claimed padding frequencies.

\subsection{Boundary drift and exact eventual record gaps}

To decide every sufficiently late crossing, extend
Lemma~\ref{lem:phase-uniform} to the seventh CDF order. Its smooth
efficiency expansion has remainder $O(m^{-4})$, including the branch
derivatives needed for optimization. The seventh-order coefficient's
value at $(p,c_*)$ will be denoted $J_*$. Its value is common to the
core comparisons and will cancel.

All derivatives below are of the explicit $Q,H$ above at $(p,c_*)$.
Put
\begin{align*}
 K(e)&=(2\alpha+Q_{pc})e+Q_p+H_c,\qquad
 F_2(e)=\alpha e^2+Q_pe+H_0,\\
 G(e)&=E_3e^3-Q_{pp}e^2/2-H_pe+K(e)^2/(2Q_{cc}).
\end{align*}
The needed constants are
\[
\begin{gathered}
 Q_{cc}=\phz(35+11s)/15,\quad
 Q_{pc}=\phz(415+131s)/540,\quad
 Q_p=\phz(1483+463s)/19440,\\
 E_3=E'''(p)/6=-\phz(19+6s)/15,\quad
 c_0'(p)=-(23+4s)/108,\quad
 Q_{pp}=Q_{\min}''(p)+Q_{cc}c_0'(p)^2.
\end{gathered}
\]
The derivatives $H_c,H_p$ are obtained by differentiating the displayed
$T_5$ expression for $H$; thus every coefficient here is explicit.
Writing $c=c_*+d_1/m$ and $v=p+e/m+d_1/m^2$, branch stationarity gives
$d_1=-K(e)/Q_{cc}$. Substituting into the branch value yields
\begin{equation}\label{eq:phase-cubic}
 B_m=\ce-Q_*/m-F_2(e)/m^2+(G(e)-J_*)/m^3+O(m^{-4})
\end{equation}
on a winning integer branch. Near a core's internal half-integer
crossing both branches must be compared. At each policy crossing in
\eqref{eq:phase-boundaries}, the two competing cores' internal phases
lie in $[T_{30},T_{02}]\subset(-1/2,1/2)$, so there is a fixed distance
from those internal crossings.

For $\theta_j=\langle\theta+jp\rangle$ and $e_j=b+\theta_j$, put
\[
 L_j(\theta)=Q_*j+F_2(e_j),\qquad
 M_j(\theta)=-Q_*j^2-2jF_2(e_j)+G(e_j).
\]
Expanding \eqref{eq:phase-cubic} at $m=N-j$ gives
\[
 B_{N-j}=\ce-Q_*/N-L_j(\theta)/N^2
                    +(M_j(\theta)-J_*)/N^3+O(N^{-4}).
\]
At a transverse limiting crossing $T$ between $i,j$, the price
difference changes sign at
\[
 \theta=T+\frac{D_{ij}}N+O(N^{-2}),\qquad
 D_{ij}=\frac{M_i(T)-M_j(T)}{L_i'(T)-L_j'(T)}.
\]
Here the equation refers to the uniform sign comparison of branch
values; no continuous interpolation of integer-threshold optima is required.
Substitution of the explicit polynomials gives
\begin{equation}\label{eq:phase-drift}
\begin{gathered}
 D_{13}=D_{02}=D_{24}=D_A
   =\frac{129918319}{529079040}
       +\frac{191267395}{1693052928}s,\\
 D_{30}=D_{41}=D_B
   =-\frac{663979357}{2351462400}
       -\frac{20072514361}{141087744000}s.
\end{gathered}
\end{equation}
Both rational parts are nonzero, so neither $sD_A$ nor $sD_B$ is rational.

The actual rotation cannot remain ambiguously close to a shifted
boundary. Let $h_N=\theta_N-T$ and $\gamma=\delta-T$.
With $L_0=25920$, the numerators of
$\gamma=(\gamma_A+\gamma_Bs)/L_0$ are
\begin{center}
\begin{tabular}{crr}
\toprule
$T$&$\gamma_A$&$\gamma_B$\\
\midrule
$T_{13}$&95720&-25573\\
$T_{30}$&72401&-18232\\
$T_{02}$&43880&-12613\\
$T_{24}$&121640&-38533\\
$T_{41}$&98321&-31192\\
\bottomrule
\end{tabular}
\end{center}
These follow directly from
$\eta/t=(60+269s)/8640$ and $3\eta/(2u)=(269+6s)/2880$.
If $k$ is the nearest integer to $Np-\delta$, then
\[
 h_N=(A_N+B_Ns)/L_0,\quad
 A_N=L_0(k-2N)+\gamma_A,\quad B_N=12960N+\gamma_B.
\]
Its field norm is
\[
 \nu_N=(A_N^2-10B_N^2)/L_0^2\in L_0^{-2}\Z.
\]
In any fixed corridor $|h_N|\le C/N$, the conjugate of $h_N$ is
$-sN+O(1)$, so $\nu_N$ is bounded and takes only finitely many rational
values. Moreover $Nh_N=-\nu_N/s+O(N^{-1})$.
For the corresponding $D_A$ or $D_B$, each $-\nu/s-D$ is nonzero by
\eqref{eq:phase-drift}. The finite set has a positive minimum separation.
Thus $|\theta_N-T-D/N|\ge\varepsilon/N$ eventually for some
$\varepsilon>0$. Outside a sufficiently wide $C/N$ corridor the same
separation is immediate. This margin dominates the $O(N^{-2})$
boundary uncertainty. Other policies have fixed limiting price gaps
near each crossing; away from crossings the unique leading winner
persists. Hence the exact sufficiently late code is the six-interval
table, shifting the three $A$ endpoints by $D_A/N$ and the two $B$
endpoints by $D_B/N$. In particular the winning core is unique.

Finally the endpoint translations
\[
 T_{02}-p=T_{13},\qquad T_{13}+q=T_{24},\qquad
 T_{30}+q=T_{41}
\]
preserve the drift group. Changing $N$ to $N+1$ changes $D/N$ only
by $O(N^{-2})$, below the norm separation. Interval transport under
$\theta_{N+1}=\langle\theta_N-p\rangle$ therefore gives exactly
\[
 0\longrightarrow1,\quad1\longrightarrow0\text{ or }2,\quad
 2\longrightarrow0\text{ or }3,\quad3\longrightarrow4,\quad
 4\longrightarrow0.
\]
The return lengths to a record are $2,3,5$. Their frequencies per
horizon are $u-\omega,t-\omega,\omega$, respectively, all positive;
this follows by translating the displayed intervals, with the vanishing
shifts leaving frequencies unchanged. Their sum is $p-\omega$, and
their gap-weighted sum is one. Each gap thus occurs infinitely often,
proving Theorem~\ref{thm:iid-padding}.
